\section{Introduction}\label{sec:intro}
Map services help travellers find routes and nearby places. Repeated trips ending near a hospital could invite inferences about a person's health or routine, even if the traveller is visiting a relative. The distinctiveness of mobility traces makes this concern relevant across repeated queries~\cite{mobilityunique}.

How much does the exchange over an entire trip help the provider distinguish
two possible trajectories? Geo-indistinguishability (Geo-I) gives a starting
point by bounding the likelihood ratio for the same output at two true
locations~\cite{geoi}. Extending that argument to a session requires accounting
even for GPS reads that help the phone decide whether to keep its previous
reference. Predictive privacy mechanisms likewise account for the cost of
fresh observations~\cite{predictive}.

For the traveller, what matters is which points of interest (POIs) appear in
the answer. Several query locations can retrieve enough candidates for the
phone to rank them using the private request~\cite{ev}. We therefore compare
the answer the user receives with the full-catalogue answer to the same request
at the true GPS position. This comparison accommodates a single replacement
coordinate, a set of real and dummy locations, or a stream with extra fake
queries. Offline trajectory releases serve a different downstream
task~\cite{transprotect,semantic,fake,sok}.

We distinguish ten scenarios by what the observer tries to infer and which
part of the exchange (the transcript) is available.

\smallskip\noindent\textbf{Location and route.}
One packet gives an estimate of the current GPS position (S1). A longer
sequence may reveal a stopping place (S2) or the route travelled, position by
position (S3). For S2 the available proxy is the mean position in the final
minute. The full session also lets us test inference of the starting point
(S9) and endpoint (S10). These are our five main localization tasks.

\smallskip\noindent\textbf{Linkage and prediction.}
Could two sessions belong to the same person or vehicle (S4)? The data
identify route families, so the test asks whether two sessions share a
family. For prediction we distinguish two decisions on the road network.
S5 asks which of two outgoing edges the user will take. S6 asks which of
two destinations the user is heading for once a turn has been observed.

\smallskip\noindent\textbf{Query content and companions.}
For S7 we hold the trip fixed and change the request on the device to see
whether its content changes the transcript. S8 considers another source of
information available to the provider. Given two time-aligned query streams,
how much does observing a companion help locate the user?


GeoI-BQ maintains a protected reference on the road network. From the reference
history the device estimates a location distribution (belief) and uses it to
choose the next query segment. We make three contributions.
\begin{itemize}
\item A common framework that specifies each inference target, the history visible to the observer and the reference used for evaluation.
\item A Geo-I-based segment policy with a session disclosure bound and conditions for service quality.
\item A six-method comparison on common tasks, with uncertainty, costs and applicable source metrics.
\end{itemize}
In the experiment the device already holds the complete static catalogue,
and the public base achieves slightly higher Recall than private adaptation.
Both findings qualify the practical benefit attributable to online retrieval.
Sections~\ref{sec:related}--\ref{sec:conclusion} cover related work, the
method and its analysis, experiments, and conclusions.
